% questions.tex —— 2012 年·辽宁卷·理科 · 题目内容片段（无答案）
% 由 tools/generate_projects.py 依来源总账自 2012/2012_数学真题_辽宁理_试题.md 生成，勿手改（改源 MD 后重跑）


\begin{gknotice}
\textbf{来源说明}：本题面依据固定源 PDF 的两页 A3 题面逐题转录，原卷印刷页码为 751—752。第 751 面为选择题 1—12、填空题 13—16 及解答题 17、18 的题头；第 752 面为解答题 19—21 与选考题 22—24。

\textbf{图形说明}：第 9 题程序框图、第 13 题三视图、第 18 题直三棱柱、第 19 题频率分布直方图、第 20 题椭圆与动圆、第 22 题两圆相切与割线，共 6 幅图由项目内 TikZ 矢量源重绘；其余各题原页无图。

\textbf{解析归属}：固定源不附答案或解析；配套答案与解析均由本库按本题面独立复算。
\end{gknotice}

\gkpart{一、选择题}

\begin{question}
\gkti{1}已知全集 $U=\{0,1,2,3,4,5,6,7,8,9\}$，集合 $A=\{0,1,3,5,8\}$，集合 $B=\{2,4,5,6,8\}$，则 $(\complement_U A)\cap(\complement_U B)=$\mbox{（　　）}

（A）$\{5,8\}$；（B）$\{7,9\}$；（C）$\{0,1,3\}$；（D）$\{2,4,6\}$。

\end{question}
\begin{question}
\gkti{2}复数 $\dfrac{2-\mathrm i}{2+\mathrm i}=$\mbox{（　　）}

（A）$\dfrac35-\dfrac45\mathrm i$；（B）$\dfrac35+\dfrac45\mathrm i$；（C）$1-\dfrac45\mathrm i$；（D）$1+\dfrac35\mathrm i$。

\end{question}
\begin{question}
\gkti{3}已知两个非零向量 $\boldsymbol a$，$\boldsymbol b$ 满足 $|\boldsymbol a+\boldsymbol b|=|\boldsymbol a-\boldsymbol b|$，则下面结论正确的是\mbox{（　　）}

（A）$\boldsymbol a\parallel\boldsymbol b$；（B）$\boldsymbol a\perp\boldsymbol b$；（C）$|\boldsymbol a|=|\boldsymbol b|$；（D）$\boldsymbol a+\boldsymbol b=\boldsymbol a-\boldsymbol b$。

\end{question}
\begin{question}
\gkti{4}已知命题 $p$：$\forall x_1,x_2\in\mathbf R$，$(f(x_2)-f(x_1))(x_2-x_1)\geqslant 0$，则 $\neg p$ 是\mbox{（　　）}

（A）$\exists x_1,x_2\in\mathbf R$，$(f(x_2)-f(x_1))(x_2-x_1)\leqslant 0$；（B）$\forall x_1,x_2\in\mathbf R$，$(f(x_2)-f(x_1))(x_2-x_1)\leqslant 0$；（C）$\exists x_1,x_2\in\mathbf R$，$(f(x_2)-f(x_1))(x_2-x_1)<0$；（D）$\forall x_1,x_2\in\mathbf R$，$(f(x_2)-f(x_1))(x_2-x_1)<0$。

\end{question}
\begin{question}
\gkti{5}一排 9 个座位坐了 3 个三口之家，若每家人坐在一起，则不同的坐法种数为\mbox{（　　）}

（A）$3\times3!$；（B）$3\times(3!)^3$；（C）$(3!)^4$；（D）$9!$。

\end{question}
\begin{question}
\gkti{6}在等差数列 $\{a_n\}$ 中，已知 $a_4+a_8=16$，则该数列前 11 项和 $S_{11}=$\mbox{（　　）}

（A）$58$；（B）$88$；（C）$143$；（D）$176$。

\end{question}
\begin{question}
\gkti{7}已知 $\sin\alpha-\cos\alpha=\sqrt2$，$\alpha\in(0,\pi)$，则 $\tan\alpha=$\mbox{（　　）}

（A）$-1$；（B）$-\dfrac{\sqrt2}{2}$；（C）$\dfrac{\sqrt2}{2}$；（D）$1$。

\end{question}
\begin{question}
\gkti{8}设变量 $x$，$y$ 满足 $\begin{cases}x-y\leqslant 10,\\0\leqslant x+y\leqslant 20,\\0\leqslant y\leqslant 15,\end{cases}$ 则 $2x+3y$ 的最大值为\mbox{（　　）}

（A）$20$；（B）$35$；（C）$45$；（D）$55$。

\Needspace{0.4\textheight}
\end{question}
\begin{question}
\gkti{9}执行如图所示的程序框图，则输出的 $S$ 值是\mbox{（　　）}
\nopagebreak[4]
\begin{center}% ============================================================
% f2012_liaoning_science_q09_flowchart.tex
% 2012 辽宁卷（理科）第 9 题：程序框图
% 来源 PDF：普通高考/2012/2012辽宁理.pdf
% 原页：第 1 页（A3），印面 751，第 9 题题干下方
% 初始 S=4、i=1；判断 i<9，是则 S=2/(2-S)、i=i+1 后循环，
% 否则输出 S 并结束。
% ============================================================
\begin{tikzpicture}[>=Stealth, line width=0.7pt, font=\small,
  every node/.style={inner sep=2pt},
  term/.style={draw, rounded corners=3pt, minimum width=2.0cm,
    minimum height=0.58cm, align=center},
  proc/.style={draw, minimum width=2.6cm, minimum height=0.58cm, align=center},
  io/.style={draw, trapezium, trapezium left angle=70, trapezium right angle=110,
    minimum width=2.4cm, minimum height=0.58cm, align=center},
  dec/.style={draw, diamond, aspect=2, minimum width=2.4cm,
    minimum height=0.9cm, align=center}]
  \node[term] (start) at (0,0) {开始};
  \node[proc] (initS) at (0,-0.92) {$S=4$};
  \node[proc] (initI) at (0,-1.84) {$i=1$};
  \node[dec] (test) at (0,-3.02) {$i<9$};
  \node[proc] (bodyS) at (0,-4.46) {$S=\dfrac{2}{2-S}$};
  \node[proc] (bodyI) at (0,-5.52) {$i=i+1$};
  \node[io] (out) at (3.35,-3.02) {输出 $S$};
  \node[term] (end) at (3.35,-4.22) {结束};

  \draw[->] (start) -- (initS);
  \draw[->] (initS) -- (initI);
  \draw[->] (initI) -- (test);
  \draw[->] (test.south) -- node[right] {是} (bodyS.north);
  \draw[->] (bodyS) -- (bodyI);
  % 循环体回到判断框左侧
  \draw[->] (bodyI.west) -- ++(-2.45,0) |- (test.west);
  % “否”分支输出
  \draw[->] (test.east) -- node[above] {否} (out.west);
  \draw[->] (out) -- (end);
\end{tikzpicture}
\end{center}
\nopagebreak[4]
（A）$-1$；（B）$\dfrac32$；（C）$\dfrac23$；（D）$4$。

\end{question}
\begin{question}
\gkti{10}在长为 $12\ \mathrm{cm}$ 的线段 $AB$ 上任取一点 $C$。现作一矩形，邻边长分别等于线段 $AC$，$CB$ 的长，则该矩形面积小于 $32\ \mathrm{cm}^2$ 的概率为\mbox{（　　）}

（A）$\dfrac16$；（B）$\dfrac13$；（C）$\dfrac23$；（D）$\dfrac45$。

\end{question}
\begin{question}
\gkti{11}设函数 $f(x)(x\in\mathbf R)$ 满足 $f(-x)=f(x)$，$f(x)=f(2-x)$，且当 $x\in[0,1]$ 时，$f(x)=x^3$。又函数 $g(x)=|x\cos(\pi x)|$，则函数 $h(x)=g(x)-f(x)$ 在 $\left[-\dfrac12,\dfrac32\right]$ 上的零点个数为\mbox{（　　）}

（A）$5$；（B）$6$；（C）$7$；（D）$8$。

\end{question}
\begin{question}
\gkti{12}若 $x\in[0,+\infty)$，则下列不等式恒成立的是\mbox{（　　）}

（A）$\mathrm e^x\leqslant 1+x+x^2$；（B）$\dfrac1{\sqrt{1+x}}\leqslant 1-\dfrac12x+\dfrac14x^2$；（C）$\cos x\geqslant 1-\dfrac12x^2$；（D）$\ln(1+x)\geqslant x-\dfrac18x^2$。

\end{question}
\gkpart{二、填空题}

\Needspace{0.4\textheight}
\begin{question}
\gkti{13}一个几何体的三视图如图所示，则该几何体的表面积为\gkblank{}。
\nopagebreak[4]
\begin{center}% ============================================================
% f2012_liaoning_science_q13_three_views.tex
% 2012 辽宁卷（理科）第 13 题：三视图
% 来源 PDF：普通高考/2012/2012辽宁理.pdf
% 原页：第 1 页（A3），印面 751，第 13 题题干下方
% 主视图 4×1，内部距左端 1、3 处有虚线；
% 侧视图 3×1，内部距左端 0.5、2.5 处有虚线；
% 俯视图 4×3 矩形，内含直径 2 的圆。
% ============================================================
\begin{tikzpicture}[>=Stealth, line width=0.7pt, font=\small,
  x=1.5cm, y=1.5cm,
  dim/.style={line width=0.4pt},
  every node/.style={inner sep=1.5pt}]
  % ---------- 主视图 ----------
  \draw (0,4) rectangle (4,5);
  \draw[dashed] (1,4) -- (1,5);
  \draw[dashed] (3,4) -- (3,5);
  % 主视图尺寸
  \draw[dim,<->] (0,3.55) -- (1,3.55) node[midway,below,font=\tiny,fill=white,inner sep=1pt] {$1$};
  \draw[dim,<->] (1,3.55) -- (3,3.55) node[midway,below,font=\tiny,fill=white,inner sep=1pt] {$2$};
  \draw[dim,<->] (3,3.55) -- (4,3.55) node[midway,below,font=\tiny,fill=white,inner sep=1pt] {$1$};

  % ---------- 侧视图 ----------
  \draw (5.3,4) rectangle (8.3,5);
  \draw[dashed] (5.8,4) -- (5.8,5);
  \draw[dashed] (7.8,4) -- (7.8,5);
  \draw[dim,<->] (5.3,3.55) -- (5.8,3.55) node[midway,below,font=\tiny,fill=white,inner sep=1pt] {$0.5$};
  \draw[dim,<->] (5.8,3.55) -- (7.8,3.55) node[midway,below,font=\tiny,fill=white,inner sep=1pt] {$2$};
  \draw[dim,<->] (7.8,3.55) -- (8.3,3.55) node[midway,below,font=\tiny,fill=white,inner sep=1pt] {$0.5$};
  % 侧视图高度
  \draw[dim,<->] (8.55,4) -- (8.55,5) node[midway,right,font=\tiny,fill=white,inner sep=1pt] {$1$};

  % ---------- 俯视图 ----------
  \draw (0,0) rectangle (4,3);
  \draw (2,1.5) circle (1);
\end{tikzpicture}
\end{center}
\nopagebreak[4]
\end{question}
\begin{question}
\gkti{14}已知等比数列 $\{a_n\}$ 为递增数列，且 $a_5^2=a_{10}$，$2(a_n+a_{n+2})=5a_{n+1}$，则数列 $\{a_n\}$ 的通项公式 $a_n=$\gkblank{}。

\end{question}
\begin{question}
\gkti{15}已知 $P$，$Q$ 为抛物线 $x^2=2y$ 上两点，点 $P$，$Q$ 的横坐标分别为 $4$，$-2$，过 $P$，$Q$ 分别作抛物线的切线，两切线交于 $A$，则点 $A$ 的纵坐标为\gkblank{}。

\end{question}
\begin{question}
\gkti{16}已知正三棱锥 $P-ABC$，点 $P$，$A$，$B$，$C$ 都在半径为 $\sqrt3$ 的球面上。若 $PA$，$PB$，$PC$ 两两互相垂直，则球心到截面 $ABC$ 的距离为\gkblank{}。

\end{question}
\gkpart{三、解答题}

\begin{question}
\gkti{17}在 $\triangle ABC$ 中，角 $A$，$B$，$C$ 的对边分别为 $a$，$b$，$c$。角 $A$，$B$，$C$ 成等差数列。

（1） 求 $\cos B$ 的值；

（2） 若边 $a$，$b$，$c$ 成等比数列，求 $\sin A\sin C$ 的值。

\Needspace{0.4\textheight}
\end{question}
\begin{question}
\gkti{18}如图，直三棱柱 $ABC-A'B'C'$，$\angle BAC=90^\circ$，$AB=AC=\lambda AA'$，点 $M$，$N$ 分别为 $A'B$ 和 $B'C'$ 的中点。
\nopagebreak[4]
\begin{center}% ============================================================
% f2012_liaoning_science_q18_prism.tex
% 2012 辽宁卷（理科）第 18 题：直三棱柱与二面角
% 来源 PDF：普通高考/2012/2012辽宁理.pdf
% 原页：第 1 页（A3），印面 751，第 18 题题干右侧
% 直三棱柱 ABC-A'B'C'，M 为 A'B 中点，N 为 B'C' 中点。
% 虚线表示被遮挡的棱 AB、AC、AA'。
% ============================================================
\begin{tikzpicture}[>=Stealth, line width=0.7pt, font=\small,
  every node/.style={inner sep=1.5pt}]
  \coordinate (A)  at (1.00,1.60);
  \coordinate (B)  at (0.20,0.20);
  \coordinate (C)  at (3.20,0.80);
  \coordinate (Ap) at (1.00,3.40);
  \coordinate (Bp) at (0.20,2.00);
  \coordinate (Cp) at (3.20,2.60);
  \coordinate (M)  at (0.60,1.80);
  \coordinate (N)  at (1.70,2.30);

  % 被遮挡的棱
  \draw[dashed] (A) -- (B);
  \draw[dashed] (A) -- (C);
  \draw[dashed] (A) -- (Ap);

  % 可见棱
  \draw (B) -- (C);
  \draw (B) -- (Bp);
  \draw (C) -- (Cp);
  \draw (Ap) -- (Bp);
  \draw (Bp) -- (Cp);
  \draw (Cp) -- (Ap);
  \draw (Ap) -- (B);

  % 中点连线与二面角相关线段
  \draw (M) -- (N);
  \draw (Ap) -- (N);
  \draw (C) -- (M);
  \draw (C) -- (N);

  % 点与标签
  \fill (A) circle (1.1pt) node[below right] {$A$};
  \fill (B) circle (1.1pt) node[below] {$B$};
  \fill (C) circle (1.1pt) node[right] {$C$};
  \fill (Ap) circle (1.1pt) node[above] {$A'$};
  \fill (Bp) circle (1.1pt) node[left] {$B'$};
  \fill (Cp) circle (1.1pt) node[right] {$C'$};
  \fill (M) circle (1.0pt) node[below left] {$M$};
  \fill (N) circle (1.0pt) node[above right] {$N$};
\end{tikzpicture}
\end{center}
\nopagebreak[4]
（1） 证明：$MN\parallel$ 平面 $A'ACC'$；

（2） 若二面角 $A'-MN-C$ 为直二面角，求 $\lambda$ 的值。

\Needspace{0.4\textheight}
\end{question}
\begin{question}
\gkti{19}电视传媒公司为了解某地区电视观众对某类体育节目的收视情况，随机抽取了 100 名观众进行调查。下面是根据调查结果绘制的观众日均收看该体育节目时间的频率分布直方图：
\nopagebreak[4]
\begin{center}% ============================================================
% f2012_liaoning_science_q19_histogram.tex
% 2012 辽宁卷（理科）第 19 题：频率分布直方图
% 来源 PDF：普通高考/2012/2012辽宁理.pdf
% 原页：第 2 页（A3），印面 752，第 19 题题干下方
% 各组频率/组距依次为 0.010、0.018、0.020、0.025、0.022、0.005，
% 组距为 10 分钟，共 6 组。
% ============================================================
\begin{tikzpicture}[>=Stealth, line width=0.7pt, font=\small,
  x=1.2cm, y=1cm,
  every node/.style={inner sep=1.5pt}]

  % 坐标轴
  \draw[->] (-0.25,0) -- (7.75,0) node[right] {分钟};
  \draw[->] (0,-0.25) -- (0,3.75) node[above] {$\dfrac{\text{频率}}{\text{组距}}$};

  % 各矩形
  \draw (0,0) rectangle (1.2,1.30);   % 0.010
  \draw (1.2,0) rectangle (2.4,2.34); % 0.018
  \draw (2.4,0) rectangle (3.6,2.60); % 0.020
  \draw (3.6,0) rectangle (4.8,3.25); % 0.025
  \draw (4.8,0) rectangle (6.0,2.86); % 0.022
  \draw (6.0,0) rectangle (7.2,0.65); % 0.005

  % 高度虚线
  \draw[dashed] (0,1.30) -- (1.2,1.30);
  \draw[dashed] (0,2.34) -- (2.4,2.34);
  \draw[dashed] (0,2.60) -- (3.6,2.60);
  \draw[dashed] (0,3.25) -- (4.8,3.25);
  \draw[dashed] (0,2.86) -- (6.0,2.86);
  \draw[dashed] (0,0.65) -- (7.2,0.65);

  % 纵轴刻度值
  \node[left,font=\scriptsize] at (0,0.65) {$0.005$};
  \node[left,font=\scriptsize] at (0,1.30) {$0.010$};
  \node[left,font=\scriptsize] at (0,2.34) {$0.018$};
  \node[left,font=\scriptsize] at (0,2.60) {$0.020$};
  \node[left,font=\scriptsize] at (0,2.86) {$0.022$};
  \node[left,font=\scriptsize] at (0,3.25) {$0.025$};

  % 横轴刻度值
  \foreach \i/\t in {0/0,1/10,2/20,3/30,4/40,5/50,6/60}
    \node[below,font=\scriptsize] at (\i*1.2,0) {$\t$};
\end{tikzpicture}
\end{center}
\nopagebreak[4]
将日均收看该体育节目时间不低于 40 分钟的观众称为“体育迷”。

（1） 根据已知条件完成下面的 $2\times2$ 列联表，并据此资料你是否认为“体育迷”与性别有关？

\begin{center}
\begin{tabular}{llll}
\toprule
 & 非体育迷 & 体育迷 & 合计 \\
\midrule
男 &  &  &  \\
女 &  & $10$ & $55$ \\
合计 &  &  &  \\
\bottomrule
\end{tabular}
\end{center}

（2） 将上述调查所得到的频率视为概率。现在从该地区大量电视观众中，采用随机抽样方法每次抽取一名观众，抽取 3 次，记被抽取的 3 名观众中的“体育迷”人数为 $X$。若每次抽取的结果是相互独立的，求 $X$ 的分布列、期望 $E(X)$ 和方差 $D(X)$。

\Needspace{10\baselineskip}
附：$\chi^2=\dfrac{n(n_{11}n_{22}-n_{12}n_{21})^2}{n_{1+}n_{2+}n_{+1}n_{+2}}$，

\begin{center}
\begin{tabular}{lll}
\toprule
$P(\chi^2\geqslant k)$ & $0.05$ & $0.01$ \\
\midrule
$k$ & $3.841$ & $6.635$ \\
\bottomrule
\end{tabular}
\end{center}

\Needspace{0.4\textheight}
\end{question}
\begin{question}
\gkti{20}如图，椭圆 $C_0:\dfrac{x^2}{a^2}+\dfrac{y^2}{b^2}=1$（$a>b>0$，$a$，$b$ 为常数），动圆 $C_1:x^2+y^2=t_1^2$，$b<t_1<a$。点 $A_1$，$A_2$ 分别为 $C_0$ 的左、右顶点，$C_1$ 与 $C_0$ 相交于 $A$，$B$，$C$，$D$ 四点。
\nopagebreak[4]
\begin{center}% ============================================================
% f2012_liaoning_science_q20_ellipse.tex
% 2012 辽宁卷（理科）第 20 题：椭圆与动圆
% 来源 PDF：普通高考/2012/2012辽宁理.pdf
% 原页：第 2 页（A3），印面 752，第 20 题题干下方
% 椭圆 C0: x^2/a^2+y^2/b^2=1，动圆 C1: x^2+y^2=t^2。
% 图中取 a=3、b=2、t=2.3 作示意，交点记为 A、B、C、D。
% ============================================================
\begin{tikzpicture}[>=Stealth, line width=0.7pt, font=\small,
  x=1.1cm, y=1.1cm,
  every node/.style={inner sep=1.5pt}]

  % 坐标轴
  \draw[->] (-3.8,0) -- (3.9,0) node[right] {$x$};
  \draw[->] (0,-2.7) -- (0,2.8) node[above] {$y$};
  \node[below left] at (0,0) {$O$};

  % 椭圆与动圆
  \draw (0,0) ellipse (3 and 2);
  \draw (0,0) circle (2.3);

  % 四个交点（a=3,b=2,t=2.3 时 x≈1.524,y≈1.722）
  \coordinate (A) at (-1.524,1.722);
  \coordinate (B) at (-1.524,-1.722);
  \coordinate (C) at (1.524,-1.722);
  \coordinate (D) at (1.524,1.722);

  \fill (A) circle (1.1pt) node[above left] {$A$};
  \fill (B) circle (1.1pt) node[below left] {$B$};
  \fill (C) circle (1.1pt) node[below right] {$C$};
  \fill (D) circle (1.1pt) node[above right] {$D$};

  % 左右顶点
  \fill (-3,0) circle (1.1pt) node[below left] {$A_1$};
  \fill (3,0) circle (1.1pt) node[below right] {$A_2$};
\end{tikzpicture}
\end{center}
\nopagebreak[4]
（1） 求直线 $AA_1$ 与直线 $A_2B$ 交点 $M$ 的轨迹方程；

（2） 设动圆 $C_2:x^2+y^2=t_2^2$ 与 $C_0$ 相交于 $A'$，$B'$，$C'$，$D'$ 四点，其中 $b<t_2<a$，$t_1\neq t_2$。若矩形 $ABCD$ 与矩形 $A'B'C'D'$ 的面积相等，证明：$t_1^2+t_2^2$ 为定值。

\end{question}
\begin{question}
\gkti{21}设 $f(x)=\ln(x+1)+\sqrt{x+1}+ax+b$（$a$，$b\in\mathbf R$，$a$，$b$ 为常数），曲线 $y=f(x)$ 与直线 $y=\dfrac32x$ 在 $(0,0)$ 点相切。

（1） 求 $a$，$b$ 的值；

（2） 证明：当 $0<x<2$ 时，$f(x)<\dfrac{9x}{x+6}$。

\end{question}
\gkpart{四、选考题（22—24 题，三选一）}

\Needspace{0.4\textheight}
\begin{question}
\gkti{22}【A】如图，$\odot O$ 和 $\odot O'$ 相交于 $A$，$B$ 两点，过 $A$ 作两圆的切线分别交两圆于 $C$，$D$ 两点，连接 $DB$ 并延长交 $\odot O$ 于点 $E$。证明：
\nopagebreak[4]
\begin{center}% ============================================================
% f2012_liaoning_science_q22_circles.tex
% 2012 辽宁卷（理科）第 22 题：两圆相交与切线
% 来源 PDF：普通高考/2012/2012辽宁理.pdf
% 原页：第 2 页（A3），印面 752，第 22 题题干下方
% 圆 O 与圆 O' 交于 A、B；过 A 的两条切线分别交两圆于 C、D；
% DB 延长交圆 O 于 E。
% 取 O(-1.2,0)、O'(1.3,0)，A(0,1.5)、B(0,-1.5)，
% 两圆半径分别为 sqrt(3.69)、sqrt(3.94)。
% ============================================================
\begin{tikzpicture}[>=Stealth, line width=0.7pt, font=\small,
  x=1.1cm, y=1.1cm,
  every node/.style={inner sep=1.5pt}]
  \coordinate (O)  at (-1.2,0);
  \coordinate (Op) at (1.3,0);
  \coordinate (A)  at (0,1.5);
  \coordinate (B)  at (0,-1.5);
  \coordinate (C)  at (-2.856,-0.975);
  \coordinate (D)  at (3.049,-0.939);
  \coordinate (E)  at (-1.788,-1.829);

  % 两圆
  \draw (O) circle ({sqrt(3.69)});
  \draw (Op) circle ({sqrt(3.94)});

  % 切线、弦与割线
  \draw (A) -- (C);
  \draw (A) -- (D);
  \draw (A) -- (B);
  \draw (D) -- (B);
  \draw (B) -- (E);
  \draw (A) -- (E);
  \draw (B) -- (C);

  % 圆心
  \fill (O) circle (1.0pt) node[below right] {$O$};
  \fill (Op) circle (1.0pt) node[below left] {$O'$};

  % 点与标签
  \fill (A) circle (1.1pt) node[above] {$A$};
  \fill (B) circle (1.1pt) node[below] {$B$};
  \fill (C) circle (1.1pt) node[below left] {$C$};
  \fill (D) circle (1.1pt) node[below right] {$D$};
  \fill (E) circle (1.1pt) node[left] {$E$};
\end{tikzpicture}
\end{center}
\nopagebreak[4]
（1） $AC\cdot BD=AD\cdot AB$；

（2） $AC=AE$。

\end{question}
\begin{question}
\gkti{23}【B】在直角坐标系 $xOy$ 中，圆 $C_1:x^2+y^2=4$，圆 $C_2:(x-2)^2+y^2=4$。

（1） 在以 $O$ 为极点，$x$ 轴非负半轴为极轴的极坐标系中，分别写出圆 $C_1$，$C_2$ 的极坐标方程，并求出圆 $C_1$，$C_2$ 的交点坐标（用极坐标表示）；

（2） 求圆 $C_1$ 与 $C_2$ 的公共弦的参数方程。

\end{question}
\begin{question}
\gkti{24}【C】已知 $f(x)=|ax+1|$（$a\in\mathbf R$），不等式 $f(x)\leqslant 3$ 的解集为 $\{x\mid -2\leqslant x\leqslant 1\}$。

（1） 求 $a$ 的值；

（2） 若 $\left|f(x)-2f\left(\dfrac x2\right)\right|\leqslant k$ 恒成立，求 $k$ 的取值范围。

\end{question}